- Fechar título.
- Troca de capítulos.
- Checar, na dedicatória, mancha do texto.
- Checar didascália do posfácio.
- Entender o que faz sentido entrar nas didascálias.
- Repensar as pretas.

---

- Rever primeiro parágrafo da apresentação.
- Fazer uma revisão dos travessões.
- Em TEXTO_PAGINAS: repensar fluidez entre linha 91 e 101.
- Em TEXTO_PAGINAS: repensar fluidez em linhas 191.
- Pensar em seções.
- Padronizar hebraico.
- Em TEXTO_PAGINAS_NOVO_CAPITULO: Dar título ao redor de BS"D. Rever início, linhas 18-45.
- Em TEXTO_CORPOS: Nota sobre "dança", linha 107.
- Em TEXTO_CORPOS: repensar fluidez em linha 137.
- Em TEXTO_CORPOS_NOVO_CAPITULO: repensar fluidez em linhas 212 a 214.
- Em TEXTO_GEOGRAFIAS: repensar fluidez em linhas 44-68.
- Em TEXTO_EPILOGO: repensar fluidez em linhas 419-421.
- Em EXTRA_POSFACIO: Há um comentário sobre a parte final no manuscrito, mas eu não entendi o que está escrito.